\section{Implementation}
\label{sec:impl}

\subsection{The interrupt only moves bytes; the main loop does the work}
At 115200 baud a byte arrives every 87~\textmu s, and the USART has a one-byte receive register, so a second byte overwrites the first
if the program is busy elsewhere (overrun). Listing~\ref{lst:irq} is the whole receive interrupt: it reads \code{RDR},
stores the byte in a 128-byte ring buffer and returns. The ring has a power-of-two size so the wrap is a mask, and it is written
only by the interrupt (\code{head}) and read only by the main loop (\code{tail}), so no lock is needed. If the ring is full the
byte is counted and dropped rather than waiting. The error flags (overrun, framing, noise) are cleared by writing \code{ICR}, as described in the reference manual~\cite{rm0091}, so that one error does not persist.
We trimmed the declarations.

\begin{figure}[t]
\begin{lstlisting}[float=false,language=C, caption={Receive interrupt and ring buffer (\texttt{uart.c}, trimmed).}, label=lst:irq]
void USART1_IRQHandler(void){
  uint32_t isr = USART1->ISR;
  if (isr & (USART_ISR_ORE|USART_ISR_FE|USART_ISR_NE)) {
    if (isr & USART_ISR_ORE) overruns++;
    USART1->ICR = USART_ICR_ORECF | USART_ICR_FECF
                | USART_ICR_NCF;    /* clear flags */
  }
  if (isr & USART_ISR_RXNE) {
    uint8_t b = (uint8_t)USART1->RDR;  /* clears RXNE */
    uint16_t next = (head + 1u) & (RING_SZ - 1u);
    if (next != tail) { ring[head] = b; head = next; }
    else dropped++;           /* full: drop, never wait */
  }
}
\end{lstlisting}
\end{figure}

\subsection{The framer replies at the 33rd byte and resynchronises at the next LF}
Listing~\ref{lst:framer} shows the byte-driven framer. The obvious alternative is to collect bytes until the LF and then test
the length, but a sender that never sends LF would then fill the buffer or stall the server. Here the 33rd byte produces
\code{NACK:TOOLONG} immediately and moves to a DISCARDING state that drops bytes silently until the next LF, so an oversize frame
gets one reply, not two, and the buffer cannot overflow. Each call writes at most one reply and never waits. The inactivity
timer is checked in \code{server_tick}, not with a delay, using the unsigned difference \code{now - last_ms}. This caused a bug: the
main loop first read the time and then drained the ring, so \code{last_ms} could be newer than \code{now} and the unsigned subtraction
wrapped to about $4\times10^9$, producing a false \code{NACK:TIMEOUT}. We found it by reading the code before the first flash and fixed
it by reading the clock \emph{after} the drain, as the comment in \code{main.c} records.

\begin{figure}[t]
\begin{lstlisting}[float=false,language=C, caption={Framer (\texttt{server.c}, trimmed). At most one reply per call.}, label=lst:framer]
size_t server_rx(server_t *s, uint8_t byte,
   uint32_t now_ms, char *out, size_t cap){
  size_t n = 0;  s->last_ms = now_ms;
  switch (s->rx) {
  case RX_IDLE: case RX_RECEIVING:
    if (byte == LF) {          /* len 0: empty frame */
      n = handle_frame(s, s->buf, s->len, out, cap);
      server_reset_framer(s);
    } else if (s->len >= SERVER_MAX_FRAME) { /*33rd*/
      n = make_reply("NACK:TOOLONG", out, cap);
      s->len = 0;  s->rx = RX_DISCARDING;
    } else { s->buf[s->len++] = (char)byte;
             s->rx = RX_RECEIVING; }
    break;
  case RX_DISCARDING:
    if (byte == LF) server_reset_framer(s);
    break;
  }
  return n;
}
\end{lstlisting}
\end{figure}

\subsection{The GUI does serial I/O on a worker thread and draws on a timer}
A \code{MOVE} can take up to the 1~s reply timeout, and a Tk window that waits for it stops redrawing and looks frozen. Listing~\ref{lst:gui}
runs each exchange on a daemon thread and passes the result back through a queue; only the Tk thread touches widgets. \code{poll}
drains the queue every 30~ms using \code{root.after}, which is Tk's thread-safe way to schedule work on its own loop. The three
exception paths separate a protocol error (\code{ReplyError}, shown as a message), and anything else such as a missing or unplugged
port (``fatal'', disconnect), so an unexpected failure can never leave the \code{busy} flag set and lock the board. The
\code{exchange} method sends \code{STATE} once after a lost reply, as described in Sec.~\ref{sec:design}.

\begin{figure}[t]
\begin{lstlisting}[float=false,language=Python, caption={GUI worker thread and queue polling (\texttt{app.py}, trimmed).}, label=lst:gui]
def run_in_thread(self, work):
    self.busy = True; self.refresh()
    def target():
        try:    self.results.put(("ok", work()))
        except ReplyError as exc:
            self.results.put(("err", str(exc)))
        except Exception as exc:  # port gone
            self.results.put(("fatal", str(exc)))
    threading.Thread(target=target, daemon=True).start()

def poll(self):
    try:
        while True:
            kind, payload = self.results.get_nowait()
            self.finish(kind, payload)
    except queue.Empty: pass
    self.root.after(30, self.poll)
\end{lstlisting}
\end{figure}
